%%% FIELD proof-prop-star-balanced-audit
\emph{Proof.}
Write \(U_j=U_j^{\mathrm{STAR}}\), \(W_j=W_j^{\mathrm{STAR}}\), and let
\(\mathcal U=\sigma(U_j:j\in\mathcal J^{\mathrm{STAR}})\).  Assumption
\(\mathrm{ass:star-complete-randomization}\) gives a conditional distribution of
\((W_j)_j\) that is constant as a function of \(\mathcal U\).  Thus, for every
assignment event \(C\) and every \(D\in\mathcal U\),
\[
 \mathbb P\{(W_j)_j\in C,\ D\}
 =\mathbb E\!\left[\mathbf 1_D
   \mathbb P\{(W_j)_j\in C\mid\mathcal U\}\right]
 =\mathbb P\{(W_j)_j\in C\}\mathbb P(D).
\]
Hence the assignment is independent of the potential-outcome vectors.  The
auxiliary randomization in Definition \(\mathrm{def:star-transformation}\) is
also independent of those vectors.  Consequently the full selected-label
array
\[
 \mathbf J=(J_{ki}:1\leq k\leq3,\ 1\leq i\leq36)
\]
is independent of \((U_j)_j\).  Its entries are pairwise distinct: within an
arm this follows from sampling without replacement, and across arms it follows
because the three assignment cells are disjoint.  The selection is feasible
because \(36=m^{\mathrm{STAR}}\leq n_k^{\mathrm{STAR}}\) for each
\(k\in\{1,2,3\}\).

For \(k\in\{1,2,3\}\), define the measurable coordinate map
\[
 g_k(y)=\left(\frac{\beta_k^{\mathrm{STAR}}y}{b_Y},0,0,0,0\right),
 \qquad -B^{\mathrm{STAR}}\leq y\leq B^{\mathrm{STAR}}.
\]
The displayed contrast satisfies
\(\max_k|\beta_k^{\mathrm{STAR}}|=1\).  Since
\(B^{\mathrm{STAR}}>0\), one has \(b_Y=B^{\mathrm{STAR}}>0\), and therefore
\[
 \|g_k(y)\|_2
 =\frac{|\beta_k^{\mathrm{STAR}}y|}{b_Y}
 \leq\frac{|\beta_k^{\mathrm{STAR}}|B^{\mathrm{STAR}}}
 {B^{\mathrm{STAR}}\max_\ell|\beta_\ell^{\mathrm{STAR}}|}
 \leq1.
\]
This verifies the asserted closed-unit-ball support, including its boundary.
Let \(A_{ki}\subseteq\mathbb B_5\) be Borel sets.  Conditional on any possible
value \(\mathbf j=(j_{ki})\) of \(\mathbf J\), Assumption
\(\mathrm{ass:star-iid-superpopulation}\), independence of \(\mathbf J\), and
distinctness of its entries give
\[
\begin{aligned}
 &\mathbb P\!\left\{X_{ki}^{\mathrm{STAR}}\in A_{ki}
       \text{ for all }k,i\mid\mathbf J=\mathbf j\right\} \\
 &\quad=
 \mathbb P\!\left\{g_k\bigl(Y_{j_{ki}}(k)\bigr)\in A_{ki}
       \text{ for all }k,i\right\} \\
 &\quad=\prod_{k=1}^3\prod_{i=1}^{36}
 Q^{\mathrm{STAR}}\!\left\{u:g_k(u_k)\in A_{ki}\right\}
 =\prod_{k=1}^3\prod_{i=1}^{36}P_k^{\mathrm{STAR}}(A_{ki}).
\end{aligned}
\]
The last expression does not depend on \(\mathbf j\), so averaging over
\(\mathbf J\) proves
\[
 \mathcal L\bigl((X_{ki}^{\mathrm{STAR}})_{k,i}\bigr)
 =\bigotimes_{k=1}^3(P_k^{\mathrm{STAR}})^{\otimes36}.
\]
This is a conditioning-on-random-labels argument; it does not require the
three auxiliary selections to be mutually independent.  The balanced
experiment contains \(3\cdot36=108\) observations, so its parameters are
\(N=108\), \(K=3\), \(m=36\), and, by construction of \(g_k\), \(d=5\).

We next give the population identification argument.  Let
\(Q_k^{\mathrm{STAR}}\) denote the marginal law of \(Y(k)\) under
\(Q^{\mathrm{STAR}}\), and let
\[
 \mathcal C^{\mathrm{STAR}}
 =\left\{\eta\in
 \mathcal P([-B^{\mathrm{STAR}},B^{\mathrm{STAR}}]^3):
 \eta_k=Q_k^{\mathrm{STAR}}\text{ for }k=1,2,3\right\}.
\]
Thus the observable distribution of the balanced data is
\(\bigotimes_k(P_k^{\mathrm{STAR}})^{\otimes36}\), whereas the unobserved
object is a coupling \(\eta\in\mathcal C^{\mathrm{STAR}}\).  For such a
coupling define the causal target
\[
 \theta(\eta)=\int
 \left(\sum_{k=1}^3\beta_k^{\mathrm{STAR}}y_k\right)^2
 \,d\eta(y_1,y_2,y_3).
\]
No cross-world restriction other than the displayed marginal and support
conditions is maintained here, so the identified set is
\[
 \Theta_I(\mathbf P^{\mathrm{STAR}})
 =\{\theta(\eta):\eta\in\mathcal C^{\mathrm{STAR}}\}.
\]

The pushforward construction now proves, rather than assumes, its relation to
the observed-law optimal-transport functional.  Put
\[
 \mathcal R_k
 =\{g_k(y):-B^{\mathrm{STAR}}\leq y\leq B^{\mathrm{STAR}}\}
 \quad\text{and}\quad
 h_k(z)=\frac{b_Yz_1}{\beta_k^{\mathrm{STAR}}}
 \quad(z\in\mathcal R_k).
\]
All three coefficients in
\(\beta^{\mathrm{STAR}}=(-1,1/2,1/2)\) are nonzero, so each inverse \(h_k\) is
well defined.  If \(G(y_1,y_2,y_3)=(g_1(y_1),g_2(y_2),g_3(y_3))\), then
\(G_\#\eta\in\Gamma(\mathbf P^{\mathrm{STAR}})\) for every
\(\eta\in\mathcal C^{\mathrm{STAR}}\).  Conversely, each
\(\gamma\in\Gamma(\mathbf P^{\mathrm{STAR}})\) assigns probability one to
\(\mathcal R_1\times\mathcal R_2\times\mathcal R_3\), and
\((h_1,h_2,h_3)_\#\gamma\in\mathcal C^{\mathrm{STAR}}\).  These two
pushforwards are inverse to one another up to null sets.  Moreover, for a
corresponding pair \((\eta,\gamma)\),
\[
\begin{aligned}
 \theta(\eta)
 &=b_Y^2\int
 \left(\sum_{k=1}^3(g_k(y_k))_1\right)^2,d\eta(y_1,y_2,y_3)\\
 &=b_Y^2\int
 \left\|\sum_{k=1}^3z_k\right\|_2^2,d\gamma(z_1,z_2,z_3),
\end{aligned}
\]
where the second equality uses the fact that every \(z_k\in\mathcal R_k\)
has four zero coordinates.  It follows that the identified set has the exact
observable-law representation
\[
 \Theta_I(\mathbf P^{\mathrm{STAR}})
 =\left\{b_Y^2\int\left\|\sum_{k=1}^3z_k\right\|_2^2,d\gamma:
 \gamma\in\Gamma(\mathbf P^{\mathrm{STAR}})\right\}.
\]

For completeness, the extrema in the last display are attained.  Here is the
direct compactness argument.  On the compact metric space
\(\mathcal X=(\mathbb B_5)^3\), take any sequence of couplings.  Choose nested
finite Borel partitions of \(\mathcal X\) whose meshes tend to zero.  Repeated
finite-dimensional subsequence extraction makes the masses of every cell in
every partition converge.  The consistent limiting cell masses define a
Borel probability measure on \(\mathcal X\).  For every continuous function
\(f\), upper and lower cell sums differ by at most the modulus of continuity
of \(f\) at the partition mesh, which tends to zero.  Thus the extracted
subsequence converges against every continuous \(f\).  Applying this first to
functions of one coordinate shows that fixed marginals are preserved, and
applying it to the continuous bounded cost
\(\|z_1+z_2+z_3\|_2^2\) shows convergence of the objective.  A minimizing and
a maximizing sequence therefore yield couplings \(\gamma_L\) and
\(\gamma_U\) attaining the two extrema.  Convex mixtures of these couplings
retain all three marginals and attain every intermediate objective value.
Consequently
\[
 \Theta_I(\mathbf P^{\mathrm{STAR}})=[\theta_L,\theta_U],
 \qquad
 \theta_L=b_Y^2\min_{\gamma\in\Gamma(\mathbf P^{\mathrm{STAR}})}
 \int\left\|\sum_{k=1}^3z_k\right\|_2^2,d\gamma,
\]
with the analogous maximum defining \(\theta_U\).  This proves sharpness, not
merely an outer bound.

Finally, Definition \(c_K(z_1,\ldots,z_K)=K^{-2}\|\sum_kz_k\|_2^2\) and
\(K=3\) imply, equation by equation,
\[
\begin{aligned}
 L(\mathbf P^{\mathrm{STAR}})
 &=\inf_{\gamma\in\Gamma(\mathbf P^{\mathrm{STAR}})}
   \int\frac{1}{9}\left\|\sum_{k=1}^3z_k\right\|_2^2,d\gamma,\\
 9L(\mathbf P^{\mathrm{STAR}})
 &=\inf_{\gamma\in\Gamma(\mathbf P^{\mathrm{STAR}})}
   \int\left\|\sum_{k=1}^3z_k\right\|_2^2,d\gamma,\\
 \theta_L&=9b_Y^2L(\mathbf P^{\mathrm{STAR}}).
\end{aligned}
\]
The first two lines also prove the unscaled equality claimed in the
proposition.  No overlap condition, matrix-rank condition, or asymptotic limit
is used; positivity of \(B^{\mathrm{STAR}}\), nonvanishing of the three
contrast coefficients, bounded support, and feasibility of the balanced
subsample are the division, inversion, support, and selection conditions used
above.  The argument establishes a guarantee for the randomized balanced
subsample of size \(108\).  The native array has cell sizes \((161,37,36)\),
not \((36,36,36)\), so no statement about a native-count plug-in rule follows
from this calculation without a separate unequal-cell theorem.  This proves
the final scope qualification. \(\square\)
